\begin{afterword}
\summaryhead

The essay distinguishes three standards. Rational evidence is anything that should change your beliefs: a police commissioner's private word that a man is a crime boss counts, since it would be foolish to insult him afterwards. Legal evidence is a narrower class. Courts will not jail someone on the commissioner's say-so, because a system that did would, over time, be abused.

Scientific knowledge is narrower again. The author's statement that he is wearing white socks is rational evidence and could be given in court, but it is not science, because nobody can test it. Science, in the essay's definition, consists of generalizations that anyone can verify by experiment without trusting authority. On this definition history is not science, and a prediction that the Sun will rise next month is, by implication, not a scientific belief either. The author says such definitions are practical choices, suggests that only openly published work should count as science, and concludes that all of it remains evidence.

\responsehead

The essay's structure is a set of three nested boxes: legal and scientific inside rational. The boxes are the argument, and the essay does not build either inner one on anything firmer than the author's assertion.

The legal box is explained by a thought experiment about a system no one proposes, jailing whoever the commissioner names, and a proverb about honey and flies. The law's own reasons are on public record. Personal knowledge and cross-examination exist to test perception, memory and narration. They are reliability rules. The essay presents them as a social convention laid over rationality. The lesson a reader takes away is that what a court refuses to hear, a rationalist may still believe on a friend's word.

The scientific box is drawn by definition, and the definition fails on contact with actual science. Science is said to be generalizations ``you can \ldots\ verify for yourself \ldots\ without having to trust anyone's authority.'' Almost no one, scientists included, holds scientific knowledge that way. The same definition expels history, and with it every science whose object happened once. The extinction of the dinosaurs is not scientific knowledge on this account. The essay's chosen historical authority, Plutarch, said he was writing lives rather than histories, and Alexander's reign and death are attested without him, by his coinage and a Babylonian astronomical diary.

The essay then strains against itself. It concedes that the sunrise prediction follows from generalizations anyone could test, and still denies it the name, with a question (``why bother to run the experiment?'') that it has already answered. When the definition is pressed, it becomes a convention: asking whether it is true is ``not a well-formed question.'' The essay keeps the authority of a definition and the immunity of a stipulation.

What the three boxes do for the reader, by inference from the essay's shape, is place the reader's own judgment in the largest one. Courts and journals impose ``special, strong, additional standards''; the rationalist is ``licensed'' by probability alone. The essay's religious verbs, ``anoint'' and ``canonize,'' mark the inner standards as ritual. The outer one is left unexamined.

In short: two of the three categories are defined by assertion, one of those definitions expels history and the historical sciences, and predictions are denied the name ``scientific'' by an argument the essay has already answered.
\end{afterword}
